\section{Common levels and positive rare mass}\label{sec:common}

We use the histories, templates and three star scores constructed in
\Cref{sec:budgets}. The first objective is to show that common levels contribute
only \(o(S)\) bit variation at a center. The second is to show that the point
exceptions on good rare levels retain a positive fraction of \(S\). Both
conclusions come from comparing the upper and lower star budgets, but the
comparison first requires a precise description of the states permitted by a
template.

Throughout this section, \(S\to\infty\), the parameters satisfy
\Cref{bud:parameters}, and \(H\) is the fixed history map from
\Cref{bud:prediction-error}. At a level, \(\nu_h\) denotes the distribution of
the true state \((A,B_1)\) of a point sampled from \(\mu\) conditional on
\(H=h\). Dependence on the level is suppressed. For a state \(y=(a,b)\), its
strict cross is
\[
 R_y=\set{(u,v):(u=a,\ v\ne b)\ \text{or}\ (u\ne a,\ v=b)}.
\]
A row means a set with fixed first coordinate; a column has fixed second
coordinate. These conventions do not refer to a geometric direction in the
plane.

\subsection{From sampled histories to the actual point}

For a center \(q\), define the mixture of predictions
\[
 \widetilde\nu_q=\E_{p\mid q}\nu_{H(p)}.
\]
For its template \(\theta=(h_1,\ldots,h_k)\), define the finite set of candidate
center states
\begin{equation}\label{com:candidates}
 \mathcal Y_\theta
 =\set{y\text{ a center state}:\nu_{h_j}(R_y)\ge1-k^{-2}
                      \text{ for every }1\le j\le k}.
\end{equation}
The actual state \(y_q\) is usually a candidate: the union bound and Markov's
inequality give, after integration on \(\mathcal C\),
\begin{equation}\label{com:actual-candidate}
 \int_{\mathcal C}\Prob(y_q\notin\mathcal Y_\theta)
 \le k^3\int_{\mathcal C}
       \E_E[1-\nu_{H(p)}(R_{y_q})]=o(S).
\end{equation}
We also need every candidate to have large mass under \(\widetilde\nu_q\).
The following elementary sampling bound supplies this uniformity.
For its first assertion, we specialize the \(\varepsilon\)-net
double-sampling argument to strict crosses
\cite[Lemmas~3.4--3.5]{HausslerWelzl1987}.
The strict-cross trace bound, the weighted-law formulation, and the subsequent
posterior-support transfer are established here.

\Needspace{5\baselineskip}
\begin{lemma}[Sampling strict crosses]\label{com:ghost}
Let \(\nu\) be a probability measure on a finite product set and let
\(X_1,\ldots,X_k\) be independent samples from \(\nu\). For
\(0<\alpha<1\) and \(k\alpha\ge8\),
\[
 \Prob\bigl(\exists y:\nu(R_y)<1-\alpha,
                    \ X_1,\ldots,X_k\in R_y\bigr)
 \le 2(2k+1)^2\,2^{-k\alpha/2}.
\]
Consequently, when \(k\) is sufficiently large, the templates in
\eqref{com:candidates} satisfy
\begin{equation}\label{com:mixture}
 \Prob\bigl(\exists y\in\mathcal Y_\theta:
              \widetilde\nu_q(R_y)<1-k^{-1/4}\mid q\bigr)
 \le 4(2k+1)^2\,2^{-k^{3/4}/2}
\end{equation}
at each level, uniformly in \(q\).
\end{lemma}

\begin{proof}
Adjoin independent samples \(X_{k+1},\ldots,X_{2k}\). If the first \(k\)
samples admit a cross as in the event, select one by a fixed ordering. For this
selected cross, the number of misses among the new samples has mean
\(k\beta\), where \(\beta=\nu(R_y^c)>\alpha\), and variance at most
\(k\beta\). Thus the probability of fewer than \(k\alpha/2\) misses is at
most \(4/(k\beta)\le1/2\). It suffices to bound twice the probability that
some cross contains the first half and misses at least \(k\alpha/2\) of the
second half.

Condition on the pooled \(2k\) samples and split their indices uniformly into
two halves. A cross induces at most \((2k+1)^2\) different subsets of these
indices: its first coordinate can agree with one of at most \(2k\) observed
coordinates, or none, and the same holds for its second coordinate. If a
pattern has \(m\) misses, the probability that they all fall in the second
half is zero for \(m>k\), and otherwise is
\[
 \frac{k(k-1)\cdots(k-m+1)}{(2k)(2k-1)\cdots(2k-m+1)}
 \le 2^{-m}.
\]
The union bound proves the first assertion. Repeated sampled states cause no
problem, since the argument concerns the indexed samples.

For the second assertion, after sampling the histories independently given
\(q\), independently sample a state from each corresponding \(\nu_{h_j}\).
Before the histories are revealed these states are independent with law
\(\widetilde\nu_q\). If a candidate violating the required conclusion exists,
select one as a function of the template. Conditional on that template, all
the synthetic states lie in its cross with probability at least
\(1-k\cdot k^{-2}=1-1/k\). Apply the first assertion with
\(\alpha=k^{-1/4}\) and divide by \(1-1/k\).
\end{proof}

Since \(W_c\le k^{1/16}S\), the integral of the failure probability in
\eqref{com:mixture} is \(o(S)\). A candidate therefore constrains the mixture
of predicted states. To constrain the actual point, we will use only rows,
columns and sets of at most two states. This restriction is essential to the
next argument.

\begin{lemma}[Transfer for restricted supports]\label{com:transfer}
Fix a level and the joint edge--template law
\[
 \Prob(p,q,\theta)=\Prob_E(p,q)\Prob(\theta\mid q).
\]
Let \(0<\eta<1/2\). As a function of \((q,\theta)\), select a support
\(U(q,\theta)\) that is a row, a column, or a set of at most two states.
On an event \(G\) determined by \((q,\theta)\), suppose
\[
 \widetilde\nu_q(U(q,\theta))\ge1-3\eta^2.
\]
Then the actual point state \(X(p)\) satisfies
\[
 \Prob\bigl(G,\ X(p)\notin U(q,\theta)\bigr)\le8\eta.
\]
The support type and its coordinates may depend on both \(q\) and \(\theta\).
\end{lemma}

\begin{proof}
Conditional on \((q,\theta)\), the history of the actual point has its
ordinary conditional law given \(q\), because the template is independent of
the point given \(q\). Hence on \(G\),
\[
 \E[\nu_{H(p)}(U(q,\theta)^c)\mid q,\theta]\le3\eta^2.
\]
The event that this support has \(\nu_{H(p)}\)-mass less than \(1-\eta\)
therefore has probability at most \(3\eta\).

For each history \(h\), define three exceptional sets using only \(\nu_h\).
If a row has mass at least \(1-\eta\), it is unique, and the first exceptional
set is its complement. If no such row exists, the first exceptional set is
empty. Define the second set in the same way for columns. Their masses are
at most \(\eta\) each.

If a set of at most two states has mass at least \(1-\eta\), put
\[
 M_h=\set{x:\nu_h(x)>\eta}
\]
and take \(M_h^c\) as the third exceptional set. Every set \(U\) of at most
two states and mass at least \(1-\eta\) contains \(M_h\): omitting an atom
of mass greater than \(\eta\) would violate that mass bound. Choosing one
such \(U\) also gives
\[
 \nu_h(M_h^c)\le\nu_h(U^c)+\nu_h(U\setminus M_h)
                 \le\eta+2\eta=3\eta.
\]
If no such \(U\) exists, the third exceptional set is empty. The union of
the three exceptional sets thus has mass at most \(5\eta\).

The actual point has marginal law \(\mu\), under which its state conditional
on \(H=h\) has distribution exactly \(\nu_h\). Consequently the probability
of this union of exceptional sets is at most \(5\eta\) under the entire
edge--template law. Outside it and the earlier nonconcentration event, the
actual state belongs to whichever permitted support was selected. Adding
the two error bounds proves the result.
\end{proof}

In the application, at most three candidates will be chosen by a fixed
ordering at each level and template. Outside the failure event bounded in
\eqref{com:mixture}, the intersection of their crosses has mixture mass at
least \(1-3k^{-1/4}\). When this intersection is contained in a row or a
column, use that row or column as \(U\); when it has size at most two, use the
intersection itself. With \(\eta=k^{-1/8}\), \Cref{com:transfer} costs at most
\[
 8k^{-1/8}W_c\le8k^{-1/16}S=o(S).
\]
An empty intersection cannot have the asserted mixture mass for large \(k\).
The proof has used the actual point marginal only for exceptional events
defined by the history and level. It has made no assertion about the law of
the actual state after additionally conditioning on a selected support.

\subsection{Classification and suppression of common variation}

We now apply this transfer to the finite geometry of strict crosses. Two
distinct center states in one row, say \((a,b)\) and \((a,b')\), have common
cross
\begin{equation}\label{com:row-intersection}
 R_{(a,b)}\cap R_{(a,b')}
       =\set{(a,v):v\notin\set{b,b'}}.
\end{equation}
A third center state outside that row cuts this intersection down to at most
one point. The corresponding assertions hold for columns. If two states
\((a,b)\) and \((a',b')\) differ in both coordinates, their common cross is
\begin{equation}\label{com:corners}
 R_{(a,b)}\cap R_{(a',b')}
       =\set{(a,b'),(a',b)}.
\end{equation}
Three states whose first coordinates are all distinct and whose second
coordinates are all distinct have empty common cross. Indeed each of the
two points in \eqref{com:corners} fails to share a coordinate with the third
state. These observations give an exhaustive classification: a candidate
set is empty, is a singleton, contains two distinct states sharing a
coordinate, or consists of states with pairwise distinct coordinates.
In the shared-coordinate case, either all candidates lie in that row or
column, or a third lies outside it. In the pairwise-distinct case there
are either two candidates or at least three.

\Cref{com:fig-crosses} shows the two-corner case: either endpoint state
together with the bit determines the state at the other endpoint.
\begin{figure}[tb]
\centering
\begin{tikzpicture}[x=1.45cm,y=1.15cm]
  \draw[gray!45] (-.55,0)--(2.55,0) (-.55,2)--(2.55,2);
  \draw[gray!45] (0,-.5)--(0,2.5) (2,-.5)--(2,2.5);
  \node[draw,circle,fill=white,inner sep=2.2pt] (y1) at (0,0) {};
  \node[draw,circle,fill=white,inner sep=2.2pt] (y2) at (2,2) {};
  \node[circle,fill=blue!65!black,inner sep=2.5pt] (x1) at (2,0) {};
  \node[circle,fill=blue!65!black,inner sep=2.5pt] (x2) at (0,2) {};
  \node[below left=3pt of y1] {\(y_1=(a,b)\)};
  \node[above right=3pt of y2] {\(y_2=(a',b')\)};
  \node[below right=3pt of x1] {\((a,b')\)};
  \node[above left=3pt of x2] {\((a',b)\)};
  \node[rotate=90,fill=white,inner sep=2pt] at (0,1) {\(Y=0\)};
  \node[rotate=90,fill=white,inner sep=2pt] at (2,1) {\(Y=0\)};
  \node[fill=white,inner sep=2pt] at (1,0) {\(Y=1\)};
  \node[fill=white,inner sep=2pt] at (1,2) {\(Y=1\)};
\end{tikzpicture}
\caption{Two candidate center states differing in both coordinates admit
exactly two point states in their common strict cross. The filled points are
those opposite corners. Within this four-state configuration, the bit of a
consistent edge and the state at either endpoint determine the state at
the other endpoint. Rows have fixed first label and columns have fixed
second label. The diagram displays coordinate
equalities of states, not geometric positions of the original plane points.}
\label{com:fig-crosses}
\end{figure}


Declare a single edge with its template bad at a common level if its state
test is inconsistent, if \(y_q\notin\mathcal Y_\theta\), if
the event inside the probability in \eqref{com:mixture} occurs, or if the
restricted-support transfer fails for
the witnesses specified by this classification. The witnesses are two
states sharing a coordinate, with a third outside that row or column when
one exists; two states differing in both coordinates; or three states with
pairwise distinct coordinates. Empty candidate sets and empty witness
intersections have no good edges. All choices are deterministic functions
of the level and template. The preceding estimates and the integrated
inconsistency bound give
\begin{equation}\label{com:bad-mass}
 \int_{\mathcal C}\Prob(p,q,\theta\text{ is bad})=o(S).
\end{equation}

The comparisons below must use the full conditional bit probability
\(\rho\) in the definition of the scores. We never remove bad edges and
renormalize that law. In particular the discontinuity of \(c(\rho)\) at
\(1/2\pm\delta\) introduces no error.

Here is the numerical estimate used in the two-corner case. For every
\(0\le\rho\le1\), with \(\delta=0.01\), \(\lambda=\delta^2/2\), and
\(c\) as in \Cref{sec:budgets},
\begin{align}
 c(\rho)\bigl(\rho^2+(1-\rho)^2\bigr)-2\rho(1-\rho)
 &= (2\rho-1)^2+
       (c(\rho)-1)\bigl(\rho^2+(1-\rho)^2\bigr)\notag\\
 &\ge\lambda/2.\label{com:score-polynomial}
\end{align}
For \(\abs{\rho-1/2}\le\delta\), discard the first square and use
\(\rho^2+(1-\rho)^2\ge1/2\). Outside that interval the expression is at
least \(4\delta^2-\lambda>\lambda/2\).

\Needspace{5\baselineskip}
\begin{proposition}[Common variation]\label{com:variation}
For the center-star and point--template-star experiments of
\Cref{bud:pair-laws},
\begin{equation}\label{com:variation-eq}
 C_Q+C_{P\theta}=o(S).
\end{equation}
Moreover their rare exception means satisfy
\begin{equation}\label{com:center-rare}
 \E_\kappa H_q^*\le\E_\mu H_p^*+o(S).
\end{equation}
\end{proposition}

\begin{proof}
Fix a common level. For each template of positive probability, let
\(\pi_\theta=\Prob(\theta)\) and
\(b_\theta=\Prob(\text{bad edge}\mid\theta)\), with both probabilities taken
under the joint law in \Cref{com:transfer}. In the center-star experiment,
condition on \(\theta\) and then sample \(q\) by its actual conditional law.
The two points are still independent with their ordinary law given \(q\).
In the point--template-star experiment, sample \(p\) conditional on
\(\theta\), then two independent centers given \((p,\theta)\). Each slot
in either experiment has the same conditional single-edge law
\(\Prob(p,q\mid\theta)\). Therefore a pair with at least one bad slot has
probability at most \(2b_\theta\) in either experiment.

Write \(V_Q(\theta)\) and \(V_{P\theta}(\theta)\) for the corresponding
averages of \(2\rho(1-\rho)\) over bases in this slice, and write
\(d_Q(\theta),d_{P\theta}(\theta)\) for its two common-level expected
scores. Each \(V\) is at most \(1/2\), and either score is at least the
negative of its variation. We prove
\begin{equation}\label{com:slice-score}
 d_Q(\theta)+d_{P\theta}(\theta)
 \ge\frac\lambda4\bigl(V_Q(\theta)+V_{P\theta}(\theta)\bigr)
       -C b_\theta
\end{equation}
with an absolute constant \(C\). All coefficients in a common score are
bounded by \(1+\lambda\), so charging a pair with a bad slot always costs
at most an absolute multiple of its probability.

\emph{An empty candidate set.}
Every edge is bad, so \(b_\theta=1\), and the boundedness of the scores proves
\eqref{com:slice-score} on increasing \(C\).

\emph{One candidate state.}
At a fixed point--template base, two good edges have the same center state.
Consistency forces their bits to agree, and the nonshared coordinate of
their center states matches. Thus this star score is at least
\(1-\lambda-O(b_\theta)\) after averaging over its bases. The center-star
score is at least \(-1/2\). Their sum is at least
\(1/2-\lambda-O(b_\theta)\), which implies \eqref{com:slice-score} because
the two variations sum to at most one.

\emph{Two candidates sharing a coordinate, and all candidates in that row
or column.}
The transfer using \eqref{com:row-intersection} places every good point in
the same row or column. If it is a row, every good edge has matching
\(A\)-coordinates and hence bit \(1\), by consistency; in the column case
every good edge has bit \(0\). A bit difference in either pair experiment
therefore requires a bad slot. This gives
\(V_Q(\theta)+V_{P\theta}(\theta)\le4b_\theta\). The lower score bound
\(d_i\ge-V_i\) now proves \eqref{com:slice-score}.

\emph{Two candidates sharing a coordinate, and a third outside that row
or column.}
The three witness crosses intersect in at most one state. If the
intersection is empty, all edges are bad. Otherwise all good point states
equal that one state. At a center base, two good edges consequently have
equal bits and the required nonshared-coordinate match. Its score is at
least \(1-\lambda-O(b_\theta)\), and the other score is at least \(-1/2\).
The same estimate as in the singleton case applies.

\emph{No coordinate shared by distinct candidates.}
If there are at least three candidates, their common cross is empty and
all edges are bad. Otherwise there are exactly two candidates. The good
point states lie in their two opposite corners from \eqref{com:corners}.
At either type of base, a good edge's bit determines the state at the other
endpoint. Thus two good edges with equal bits have the required collision.
For the full conditional probability \(\rho\) at that base, their expected
score is at least the left side of \eqref{com:score-polynomial}, less an
absolute multiple of the probability of a bad slot. After averaging bases,
each score is at least \(\lambda/2-O(b_\theta)\), which proves
\eqref{com:slice-score}.

This completes the classification. Multiply \eqref{com:slice-score} by its
actual slice weight \(\pi_\theta\), sum, and integrate over common levels.
The error is
\(C\int_{\mathcal C}\sum_\theta\pi_\theta b_\theta=o(S)\) by
\eqref{com:bad-mass}. This argument allows individual \(b_\theta\) to be
large; no inverse slice weight or number of templates enters. The common
contributions to \(D_Q+D_{P\theta}\) are therefore at least
\[
 \frac\lambda4(C_Q+C_{P\theta})-o(S).
\]
The lower bounds in \Cref{bud:rare-lower} outside common levels cancel their
main terms, because the two star vertices are on opposite sides. Hence
\[
 \frac\lambda4(C_Q+C_{P\theta})-o(S)
 \le D_Q+D_{P\theta}
 \le 10\delta^2\lambda(C_Q+C_{P\theta})+o(S),
\]
where the last inequality is \Cref{bud:upper}. Since
\(10\delta^2<1/4\), this proves \eqref{com:variation-eq}.

Finally the common contribution to \(D_Q\) is at least \(-C_Q\).
\Cref{bud:rare-lower} and \Cref{bud:upper} therefore give
\[
 2\E_\kappa H_q^*-2\E_\mu H_p^*-C_Q-o(S)
 \le D_Q\le10\delta^2\lambda C_Q+o(S).
\]
Use \(C_Q=o(S)\) to obtain \eqref{com:center-rare}.
\end{proof}

\subsection{A second score forces positive rare mass}

We have suppressed common variation at centers. To show that point exceptions
remain on rare levels, suppose temporarily that their mean is \(o(S)\).
\Cref{com:variation} then gives the same bound for center exceptions. We will
show that common variation is also small at ordinary point stars. Once both
ordinary star variations are small, expansion forces the edge bit itself to
be almost constant at nearly every level, contradicting \(S=\int f\).

For a probability measure \(\nu\) on states, let \(\nu_A,\nu_{B_1}\) denote
its coordinate marginals, with squared norms equal to the sums of squares of
their atom masses. Define
\begin{equation}\label{com:T}
 T=\int_{\mathcal C}\E_\mu\left[
       \norm{(\nu_{H(p)})_A}_2^2+
       \norm{(\nu_{H(p)})_{B_1}}_2^2-1\right].
\end{equation}
These norms count coordinate matches between two independent predicted
points with the same history. Unlike the preceding star scores, \(T\) uses
no center.

\begin{lemma}[The two-prediction score]\label{com:prediction-score}
If \(\E_\mu H_p^*=o(S)\), then \(T\le o(S)\).
\end{lemma}

\begin{proof}
Choose a history with its marginal probability, then two independent points
\(p',p''\) from \(\mu\) conditional on that history. This is one fixed
coupling throughout all levels, and each point marginal is \(\mu\). At a
level, the expected value of
\[
 F(p',p'')=
   \1_{\{A(p')=A(p'')\}}+
   \1_{\{B_1(p')=B_1(p'')\}}-1
\]
is the integrand in \eqref{com:T}.

For a distinct pair put \(Z=z_{p'}-z_{p''}\) and
\(W=w_{p'}-w_{p''}\), both nonzero. At a place with level interval
\([l,u]\), the endpoint conditions from \eqref{ht:range} say
\(l\le-\log|Z|_v\) and \(u\ge\log|W|_v\). The \(A\)-match can hold only
for \(s\le-\log|Z|_v+O_v(1)\), and the \(B_1\)-match only for
\(s\ge\log|W|_v-O_v(1)\). Here the constants vanish at nonarchimedean
places and are absolute at infinity. Therefore
\[
 \int_l^u F(p',p'')\dd s
 \le-\log|Z|_v-\log|W|_v+O_v(1).
\]
The product formula, and total archimedean weight one, show that the
integral over all levels is at most an absolute constant.

At a good rare level where neither point is exceptional, their dominant
coordinate labels match, so \(F\ge0\). In general at such a level
\[
 F\ge-\1_{\{p'\text{ exceptional}\}}-
         \1_{\{p''\text{ exceptional}\}}.
\]
The contribution on damaged levels is at least \(-W_b\). Removing all
noncommon levels from the full integral consequently gives, for each
distinct pair,
\[
 \int_{\mathcal C}F(p',p'')
       \le O(1)+H_{p'}^*+H_{p''}^*+W_b.
\]

On the diagonal the common-level integral is exactly \(W_c\). Its
probability under the selected coupling is
\[
 \sum_{p\in P}\frac{\mu(p)^2}{\mu(H^{-1}(H(p)))}
 \le |H(P)|\max_p\mu(p)=t^{-1+o(1)}.
\]
Here \(|H(P)|\le L^{2K}=t^{o(1)}\) and the marginal degree bounds were
proved in \Cref{sec:incidence}. Since
\(W_c\le k^{1/16}S\) and \(k=t^{o(1)}\), its contribution is \(o(S)\).
Average the bound for distinct pairs, using the two marginals \(\mu\),
\(W_b=o(S)\), and \(S\to\infty\). This proves
\(T\le2\E_\mu H_p^*+o(S)=o(S)\).
\end{proof}

\begin{proposition}[Positive rare mass]\label{com:rare-positive}
The good rare point exceptions satisfy
\begin{equation}\label{com:rare-positive-eq}
 \liminf_{t\to\infty}\frac{\E_\mu H_p^*}{S}>0.
\end{equation}
\end{proposition}

\begin{proof}
If the conclusion fails, restrict to a subsequence on which
\(\E_\mu H_p^*=o(S)\). By \eqref{com:center-rare}, also
\(\E_\kappa H_q^*=o(S)\), and \Cref{com:prediction-score} gives
\(T\le o(S)\).

At a common level declare an ordinary edge bad if its test is inconsistent
or if
\(\nu_{H(p)}(R_{y_q})<1-k^{-2}\). Its integrated probability is \(o(S)\),
because Markov's inequality bounds the second contribution by
\(k^2\int_{\mathcal C}\E_E[1-\nu_{H(p)}(R_{y_q})]=o(S)\).
Fix a point \(p\) and this level. Write \(b_p\) for its conditional bad-edge
probability, \(\nu=\nu_{H(p)}\), \(\zeta=k^{-2}\), and
\[
 \tau(\nu)=\norm{\nu_A}_2^2+\norm{\nu_{B_1}}_2^2-1.
\]
Let \(\mathcal Y_p\) be the states of its good neighbors. Every such state
has \(\nu(R_y)\ge1-\zeta\), and the actual point state belongs to each
\(R_y\), by consistency. Let \(d_P(p)\) be the common-level score at this
ordinary point star, using its full bit probability \(\rho_p\). We claim
\begin{equation}\label{com:second-slice-score}
 d_P(p)+\tau(\nu)
 \ge\frac\lambda4\,2\rho_p(1-\rho_p)-C(b_p+\zeta).
\end{equation}
As before, a pair with a bad slot costs only \(O(b_p)\); its probability is
at most \(2b_p\). The strict-cross classification proves the claim as
follows.

If \(\mathcal Y_p\) is empty, \(b_p=1\), while \(d_P(p)\ge-1/2\) and
\(\tau(\nu)\ge-1\), so the claim holds for an absolute \(C\).

If \(\mathcal Y_p=\set{(a,b)}\), put
\(r=\nu_A(a)\), \(s=\nu_{B_1}(b)\). Since
\(r+s\ge\nu(R_{(a,b)})\ge1-\zeta\),
\[
 \tau(\nu)\ge r^2+s^2-1
        \ge\tfrac12(1-\zeta)^2-1\ge-\tfrac12-\zeta.
\]
Two good neighbors have the same state and bit, so
\(d_P(p)\ge1-\lambda-O(b_p)\). Their sum has the positive lower bound
needed for \eqref{com:second-slice-score}.

If two states share a row, their common cross is in that row, of
\(\nu\)-mass at least \(1-2\zeta\). Thus
\(\tau(\nu)\ge(1-2\zeta)^2-1\ge-4\zeta\).
If every state is in that row, the actual point lies there too, since it
belongs to the crosses of both distinct states. All good edges consequently
have bit \(1\). Their full-law variation is at most \(2b_p\), and
\(d_P(p)\ge-2b_p\), proving the claim in this case. Columns give bit \(0\)
and the same estimates. If there is a third state outside the row or column,
the three crosses have intersection of size at most one and mass at least
\(1-3\zeta\). It must be a singleton for large \(k\). Both marginals then
have an atom of at least that mass, giving
\[
 \tau(\nu)\ge2(1-3\zeta)^2-1\ge1-12\zeta.
\]
Together with \(d_P(p)\ge-1/2\), this is more than sufficient.

Finally suppose no coordinate is shared by distinct states in
\(\mathcal Y_p\). Three such states are impossible, since the actual point
belongs to all their crosses and their intersection is empty. In the
remaining case there are two states differing in both coordinates. The two
opposite corners carry \(\nu\)-mass at least \(1-2\zeta\). Each coordinate
marginal therefore has at least that mass on a set of size two, so
\[
 \tau(\nu)\ge
  2\cdot\frac{(1-2\zeta)^2}{2}-1\ge-4\zeta.
\]
The actual point is one of those corners. Its bit determines the good
neighbor state; hence \eqref{com:score-polynomial} gives
\(d_P(p)\ge\lambda/2-O(b_p)\). This completes the proof of
\eqref{com:second-slice-score}.

Average \eqref{com:second-slice-score} over \(p\) and integrate on common
levels. The bad-edge integral is \(o(S)\), and
\(\zeta W_c\le k^{-31/16}S=o(S)\). The result is
\[
 (\text{common contribution to }D_P)+T
       \ge\frac\lambda4 C_P-o(S).
\]
The noncommon contribution to \(D_P\) is at least \(-o(S)\) by
\Cref{bud:rare-lower}, under the present assumptions on both exception
means. Therefore
\[
 \frac\lambda4 C_P-o(S)
 \le D_P+T\le10\delta^2\lambda C_P+o(S),
\]
and again \(10\delta^2<1/4\) yields \(C_P=o(S)\).

To finish, let \(V_P,V_Q\) be the ordinary star variations at any fixed
level. Assign to each vertex its conditional majority bit, breaking ties
in a fixed way. The probability that the bit on an incident edge differs
from this label is at most the corresponding average variation, because
\(\min(\rho,1-\rho)\le2\rho(1-\rho)\). Thus the two endpoint labels disagree
on an edge with probability at most \(V_P+V_Q\). When that sum is small,
the matching-label consequence of expansion in \Cref{inc:matching} gives
one bit label with total marginal exception probability
\(O(V_P+V_Q)\). The edge bit differs from it with probability
\(O(V_P+V_Q)\) as well. When the sum is not small the same bound follows
from \(f\le1/2\), with a larger absolute constant. Consequently
\[
 f\le C(V_P+V_Q)
\]
uniformly for all sufficiently large \(t\). Integrating on common levels
gives \(\int_{\mathcal C}f\le C(C_P+C_Q)=o(S)\).

On good rare levels, the minority-bit probability is bounded by the two
marginal exception probabilities and the inconsistency probability, as in
\Cref{bud:rare-moments}. Its integral is \(o(S)\) under the present
assumptions. Damaged levels contribute at most \(W_b/2=o(S)\). Thus
\(S=\int f=o(S)\), a contradiction. This proves
\eqref{com:rare-positive-eq}.
\end{proof}

The positive rare mass is a first-moment statement. To extract many
pairs at a single height scale, \Cref{sec:extraction} next controls its
large-height tail and removes errors in height-weighted measure before
selecting a band. Throughout that selection, the center-star variation
bound $C_Q=o(S)$ remains available under the original wedge law.
